Запросы на выполнение команд пульта РК/ПКУ отправляются из вкладки <<РК/ПКУ>> основного окна программы. Существуют следующие команды:

\begin{itemize}
	\item получить версию драйвера;
	\item проверить соединение;
	\item проверить работоспособность пульта;
	\item прочитать основную информацию пульта;
	\item записать основную информацию пульта;
	\item прочитать длительности ПКУ;
	\item установить режим приёма одной ПКУ;
	\item выдать одну РК.
\end{itemize}

Команда получения версии драйвера вызывается с помощью кнопки <<Проверить версию>>. Версия отображается под кнопкой.

Команда проверки соединения вызывается с помощью кнопки <<Проверить соединение>>. Результат проверки отображается под кнопкой.

Команда проверки работоспособности пульта вызывается с помощью кнопки <<Статус оборудования>>. Результат проверки отображается под кнопкой.

Команда чтения основной информации пульта вызывается с помощью кнопки <<Читать основную информацию>>. Надпись под кнопкой показывает, успешно ли чтение. Основная информация отображается в нижнем текстовом поле. Перечисляются следующие поля основной информации:

\begin{itemize}
	\item идентификатор модуля в кодировке UTF-8;
	\item размер буфера на приём;
	\item описание модуля в кодировке UTF-8;
	\item MAC-адрес в виде <<$\mathtt{AA}:\mathtt{BB}:\mathtt{CC}:\mathtt{DD}:\mathtt{EE}:\mathtt{FF}$>>, где между двоеточиями указаны байты в шестнадцатеричной системе счисления;
	\item IP-адрес в виде <<$\mathtt{a.b.c.d}$>>, где между точками указаны байты в десятичной системе счисления;
	\item маска подсети в виде <<$\mathtt{a.b.c.d}$>>, где между точками указаны байты в десятичной системе счисления;
	\item основной шлюз в виде <<$\mathtt{a.b.c.d}$>>, где между точками указаны байты в десятичной системе счисления;
	\item адрес DNS-сервера в виде <<$\mathtt{a.b.c.d}$>>, где между точками указаны байты в десятичной системе счисления;
	\item флаг <<Использовать DHCP>> ($0$ или $1$).
\end{itemize}

Команда записи основной информации вызывается с помощью кнопки <<Записать основную информацию>>. Для записи необходимо заполнить все поля на панели <<Параметры основной информации>>:

\begin{itemize}
	\item описание вводится как не более чем $120$ символов в кодировке UTF-8;
	\item MAC-адрес вводится в виде <<$\mathtt{AA:BB:CC:DD:EE:FF}$>>, где между двоеточиями указаны байты в шестнадцатеричной системе счисления;
	\item IP-адрес вводится в виде <<$\mathtt{a.b.c.d}$>>, где между точками указаны байты в десятичной системе счисления;
	\item маска подсети вводится в виде <<$\mathtt{a.b.c.d}$>>, где между точками указаны байты в десятичной системе счисления;
	\item основной шлюз вводится в виде <<$\mathtt{a.b.c.d}$>>, где между точками указаны байты в десятичной системе счисления;
	\item адрес DNS-сервера вводится в виде <<$\mathtt{a.b.c.d}$>>, где между точками указаны байты в десятичной системе счисления;
	\item флаг <<Использовать DHCP>> устанавливается и сбрасывается с помощью флажка <<Использовать DHCP>>.
\end{itemize}

Надпись под кнопкой <<Записать основную информацию>> показывает, успешна ли запись.

Команда чтения длительностей ПКУ вызывается с помощью кнопки <<Читать длительности ПКУ>>. Для чтения необходимо выбрать номера ПКУ в списке <<Выбор номеров ПКУ>>. Надпись под списком показывает, успешно ли чтение. Прочитанные длительности отображаются в нижнем текстовом поле в миллисекундах.

Команда установки режима приёма одной ПКУ вызывается с помощью кнопки <<Применить режим>>. Для установки нужно выбрать номер ПКУ в выпадающем списке <<Выбор ПКУ>> и режим в выпадающем списке <<Отрицательный импульс>>. Надпись под выпадающим списком <<Выбор ПКУ>> показывает, успешно ли выполнение команды.

Команда выдачи одной РК вызывается с помощью кнопки <<Выдать РК>>. Для выдачи нужно заполнить поля ввода. В поле <<Номер РК>> указывается номер РК от $1$ от $48$ включительно. В поле <<Длительность (мс)>> указывается длительность РК в миллисекундах как целое неотрицательное число. Надпись под полем <<Номер РК>> показывает, успешно ли выполнение команды.